\subsection{Entrelacement APP--Witt, voisin de Leech et orbifold d'ordre \(3\)}

Les douze fibres portent les étiquettes
\[
 (i,\mu,\varepsilon)\in
 \mathbb Z/3\mathbb Z\times\{0,1\}\times\{0,1\},
 \qquad
 \lambda(i,\mu,\varepsilon)=4i+2\mu+\varepsilon\pmod{12}.
\]
L'application \(\lambda\) est bijective et, les trivialisations et les cadres étant fixés,
les entrelaceurs satisfont
\[
 g_{\lambda(i,\mu,\varepsilon)}T_{i\mu\varepsilon}
 =T_{i\mu\varepsilon}C^{(-1)^\mu}.
\]
Leur somme directe est un isomorphisme \(C_3\)-équivariant
\[
 T_\lambda:W_{\mathrm{APP}}\otimes\mathbb C
 \xrightarrow{\ \cong\ }
 \bigoplus_{b=0}^{11}(A_{2,b}\otimes\mathbb C).
\]

Soit \(N=N(A_2^{12})\), soit \(v=v_\alpha\) le vecteur radial de norme \(54\),
et définissons
\[
 N_0=\ker\bigl(\langle\,\cdot\,,v\rangle\bmod3\bigr),
 \qquad N_v=N_0+\mathbb Z\frac v3.
\]
Le calcul réticulaire vérifie que \(N_v\cong\Lambda_{24}\), que l'
isométrie \(g_c\) est d'ordre trois et ne possède aucun point fixe non nul, et que
\[
 \frac{g_cv-v}{3},\quad\frac{g_c^{-1}v-v}{3}\in N_0.
\]
Par conséquent, \(g_cN_v=N_v\). Son relèvement standard au VOA réticulaire possède un poids
tordu \(4/3\) et un type zéro. Les théorèmes classiques d'orbifold cyclique donnent
\[
 V_{(3)}:=\operatorname{Orb}_{\mathbb Z/3}
       (V_{N_v},\widehat g_c)\cong V^\natural,
 \qquad T_{3B}=t_3+12.
\]
Le caractère se retrouve par la moyenne des neuf secteurs et par l'
uniformisation de niveau trois :
\[
 \operatorname{ch}_{V_{(3)}}
 =\frac13\sum_{i,j\in\mathbb Z/3}Z_{i,j}=J,
\]
\[
 J=t_3+12+\frac{10\,3^9}{t_3}
       +\frac{4\,3^{14}}{t_3^2}+\frac{3^{18}}{t_3^3}.
\]
En particulier,
\[
 [q]J=54+10\cdot3^9=196884=54(5\cdot729+1).
\]

\subsection{Constructions d'ordres \(2\) et \(3\) de l'algèbre d'opérateurs de vertex \texorpdfstring{\(V^\natural\)}{V natural}}

Fixons l'extension centrale et le cocycle de Leech
\[
 1\longrightarrow\{\pm1\}\longrightarrow\widehat\Lambda_{24}
 \longrightarrow\Lambda_{24}\longrightarrow1,
 \qquad
 V_\Lambda=M(1)\otimes\mathbb C_\varepsilon[\Lambda_{24}].
\]
La négation \(\lambda\mapsto-\lambda\) se relève en \(\theta\), et
\[
 t_2(\tau)=\operatorname{Tr}_{V_\Lambda}
  \left(\theta q^{L_0-1}\right)
 =q^{-1}\prod_{n\ge1}(1+q^n)^{-24}.
\]
La présentation FLM est
\[
 V_{(2)}:=(V_\Lambda)^+\oplus(V_\Lambda^T)^+\cong V^\natural.
\]

\begin{theorem}[Équivalence typée des présentations d'ordres \(3\) et \(2\)]
\label{thm:cierre-vtm-equivalencia-voa-23}
Soient
\[
 \phi_3:V_{(3)}\xrightarrow{\ \cong\ }V^\natural,
 \qquad
 \phi_2:V_{(2)}\xrightarrow{\ \cong\ }V^\natural
\]
les isomorphismes gradués fournis par les théorèmes d'orbifold
appliqués aux données précédentes. Alors
\[
 \boxed{\Phi_{23}:=\phi_2^{-1}\circ\phi_3:
 V_{(3)}\xrightarrow{\ \cong\ }V_{(2)}}
\]
est un isomorphisme gradué d'algèbres d'opérateurs de vertex. De plus,
\[
 a\longmapsto\Phi_{23}a\Phi_{23}^{-1}
\]
identifie leurs groupes d'automorphismes et conserve les traces graduées.
\end{theorem}

\begin{proof}
La composition d'isomorphismes de VOA est un isomorphisme de VOA et préserve le
vide, le vecteur conforme, les opérateurs de vertex et la graduation par
\(L_0\). La conjugaison transporte les automorphismes. Pour chaque automorphisme
\(a\), l'invariance de la trace par conjugaison donne l'égalité des
séries graduées correspondantes.
\end{proof}

L'isomorphisme \(\Phi_{23}\) dépend des identifications \(\phi_2\) et
\(\phi_3\) ; aucune canonicité absolue n'est affirmée. La construction d'ordre deux
conserve en outre le contenu exclusif
\[
 T_{2A}=t_2+24+\frac{2^{12}}{t_2},
 \qquad
 \boxed{j=\frac{(t_2+256)^3}{t_2^2}},
\]
qui complète l'uniformisation modulaire d'ordre deux.

Une fois \(V^\natural\) construit, les trois publications s'énoncent sans
les sérialiser indûment :
\[
 \operatorname{Aut}(V^\natural)\cong\mathbb M,
 \qquad
 \operatorname{ch}_{V^\natural}=J,
 \qquad
 g\longmapsto T_g(\tau)
 =\operatorname{Tr}_{V^\natural}\!\left(gq^{L_0-1}\right).
\]
Le théorème de Moonshine coordonne l'action graduée avec les fonctions de
genre zéro et leurs identités de réplicabilité. L'action du Monstre vit
dans \(V^\natural\) ; le fait qu'elle ne soit pas une permutation ponctuelle de chiffres est un contrôle
final des types, non la thèse du corridor.

